\begin{appendix}

\subsection*{附录A \quad 端到端加密核心代码}

\subsubsection*{A.1 \quad SealedMessage 信封实现（Kotlin）}

\begin{lstlisting}[language=Java]
object SealedMessage {
    private const val VERSION = 1
    private const val NONCE_BYTES = 12
    private const val GCM_TAG_BITS = 128

    fun seal(mkey: ByteArray, plaintext: String): String {
        val nonce = ByteArray(NONCE_BYTES)
        SecureRandom().nextBytes(nonce)
        val cipher = Cipher.getInstance("AES/GCM/NoPadding")
        val spec = GCMParameterSpec(GCM_TAG_BITS, nonce)
        cipher.init(Cipher.ENCRYPT_MODE, SecretKeySpec(mkey, "AES"), spec)
        val ct = cipher.doFinal(plaintext.toByteArray(Charsets.UTF_8))
        val json = """{"v":$VERSION,
            "n":"${b64uEncode(nonce)}",
            "c":"${b64uEncode(ct)}"}"""
        return b64uEncode(json.toByteArray(Charsets.UTF_8))
    }

    fun tryOpen(mkey: ByteArray, sealed: String): String? {
        return try {
            val outer = String(b64uDecode(sealed), Charsets.UTF_8)
            val obj = JSONObject(outer)
            if (obj.optInt("v") != VERSION) return null
            val nonce = b64uDecode(obj.getString("n"))
            val ct = b64uDecode(obj.getString("c"))
            val cipher = Cipher.getInstance("AES/GCM/NoPadding")
            val spec = GCMParameterSpec(GCM_TAG_BITS, nonce)
            cipher.init(Cipher.DECRYPT_MODE, SecretKeySpec(mkey, "AES"), spec)
            String(cipher.doFinal(ct), Charsets.UTF_8)
        } catch (_: Exception) { null }
    }
}
\end{lstlisting}

\subsubsection*{A.2 \quad 密钥派生实现（Kotlin）}

\begin{lstlisting}[language=Java]
object CryptoUtils {
    private const val SEED_SALT = "stegochat-seed-v1"
    private const val MKEY_SALT = "stegochat-mkey-v1"
    private const val MKEY_INFO = "e2ee"
    private const val PBKDF2_ITER = 100_000
    private const val USER_KEY_BITS = 256
    private const val MKEY_BYTES = 32
    private const val FINGERPRINT_BYTES = 8

    fun deriveUserKey(phrase: String): ByteArray {
        val factory = SecretKeyFactory.getInstance("PBKDF2WithHmacSHA256")
        val spec = PBEKeySpec(
            phrase.toCharArray(),
            SEED_SALT.toByteArray(Charsets.UTF_8),
            PBKDF2_ITER, USER_KEY_BITS
        )
        return factory.generateSecret(spec).encoded
    }

    fun deriveMkey(selfKey: ByteArray, peerKey: ByteArray): ByteArray {
        val (lo, hi) = if (selfKey.toHexString() <= peerKey.toHexString())
            selfKey to peerKey else peerKey to selfKey
        val ikm = lo + hi
        return hkdfSha256(
            ikm = ikm,
            salt = MKEY_SALT.toByteArray(Charsets.UTF_8),
            info = MKEY_INFO.toByteArray(Charsets.UTF_8),
            length = MKEY_BYTES
        )
    }

    fun fingerprint(userKey: ByteArray): String {
        val hash = MessageDigest.getInstance("SHA-256").digest(userKey)
        return hash.take(FINGERPRINT_BYTES)
            .joinToString(":") { "%02x".format(it) }
    }
}
\end{lstlisting}

\subsection*{附录B \quad WebSocket 连接管理核心代码}

\subsubsection*{B.1 \quad 服务端 ConnectionManager}

\begin{lstlisting}[language=Python]
class ConnectionManager:
    def __init__(self):
        self.active: dict[str, WebSocket] = {}

    async def connect(self, user_id: str, ws: WebSocket):
        await ws.accept()
        old = self.active.get(user_id)
        if old:
            try:
                await old.send_text(json.dumps({
                    "type": "kicked",
                    "payload": {
                        "reason": "Logged in on another device"
                    }
                }))
            except Exception:
                pass
            async def _force_close():
                await asyncio.sleep(5)
                try:
                    await old.close(code=4001)
                except Exception:
                    pass
            asyncio.create_task(_force_close())
        self.active[user_id] = ws

    async def send_to_user(self, user_id: str,
                           message: dict) -> bool:
        ws = self.active.get(user_id)
        if ws:
            try:
                await ws.send_text(json.dumps(message))
                return True
            except Exception:
                self.disconnect(user_id, ws)
        return False
\end{lstlisting}

\subsection*{附录C \quad SparSamp 核心算法代码}

\subsubsection*{C.1 \quad 消息驱动采样与稀疏采样嵌入}

\begin{lstlisting}[language=Python]
def gaussian_quantize_256(mu, sigma, low=-1.0, high=1.0):
    """将 N(mu, sigma) 量化为 256 个 bin 的概率分布"""
    edges = np.linspace(low, high, 257)
    cdf = norm.cdf(edges, loc=mu, scale=sigma)
    probs = np.diff(cdf)
    probs = np.clip(probs, 1e-30, None)
    return probs / probs.sum()

def encode_step(bin_idx, probs, n_m, k_m, block_size):
    """单像素消息驱动采样步骤"""
    r = random.random()
    cum = np.cumsum(probs)
    r_i_m = (k_m / n_m + r) % 1.0
    selected = int(np.searchsorted(cum, r_i_m))
    selected = min(selected, len(probs) - 1)

    cum_full = np.concatenate(([0.0], cum))
    lo = cum_full[selected]
    hi = cum_full[selected + 1]
    n_m_new = int(np.floor(n_m * (hi - lo)))
    k_m_new = int(np.floor(k_m - n_m * lo))
    n_m_new = max(n_m_new, 1)

    done = (n_m_new <= 1)
    return selected, n_m_new, k_m_new, done
\end{lstlisting}

\end{appendix}
